%=====================================================================
\chapter{Natural Language Processing}\label{ch:nlp}
%=====================================================================
\locator{6}{Part VI --- Learning, Language and Responsible AI}

\begin{outcomes}
By the end of this chapter you will be able to:
\begin{tight}
  \item \textbf{Describe} the symbolic language pipeline from characters to
        meaning, and \textbf{say} what each stage contributes.
  \item \textbf{Write} a context-free grammar for a small sub-language and
        \textbf{parse} sentences with the CYK algorithm.
  \item \textbf{Identify} structural ambiguity, \textbf{count} the parses, and
        \textbf{explain} why the ambiguity is operationally consequential.
  \item \textbf{Implement} pattern substitution in the style of Eliza, and
        \textbf{say} precisely what it demonstrates and what it does not.
  \item \textbf{Build} an n-gram model and \textbf{state} the limitation that no
        amount of data removes.
\end{tight}
\end{outcomes}

\begin{prereq}
\chref{python} (dictionaries, recursion) and \chref{fol} --- pattern matching with
variables is unification in another guise, and recognising that saves learning it
twice. \chref{uncertainty} for the n-gram section.
\end{prereq}

\begin{keyterms}
tokenisation $\bullet$ morphology $\bullet$ stemming $\bullet$ part of speech
$\bullet$ context-free grammar $\bullet$ terminal and non-terminal $\bullet$
Chomsky normal form $\bullet$ parse tree $\bullet$ chart parsing $\bullet$ CYK
$\bullet$ structural ambiguity $\bullet$ attachment $\bullet$ pattern substitution
$\bullet$ Eliza $\bullet$ n-gram $\bullet$ Markov assumption $\bullet$ sparsity
\end{keyterms}

\section{Ticket Subject Lines}

The service desk receives several hundred tickets a day, each with a one-line
subject written by a person in a hurry:

\begin{quote}\ttfamily\small
restart the server after the timeout\\
check the payments gateway\\
fix the login error on the gateway
\end{quote}

The routing system needs to know what is being asked for and what it is being asked
about. That is a language problem, and this chapter takes the \emph{symbolic}
approach to it: grammars, parsing and pattern matching, with no learned component
anywhere.

That approach is not how a modern system would do it, and the reasons to study it
are specific. Grammars are the clearest case of a representation whose structure
\emph{is} the knowledge. Ambiguity is easiest to see when a parser hands you two
trees rather than one confident answer. And Eliza, in Section~\ref{sec:eliza}, is
the most instructive failure in the history of the field --- it is impossible to
understand what modern language systems are and are not doing without it.

\section{The Symbolic Pipeline}

\dsfig{%
\begin{tikzpicture}[font=\scriptsize,
  st/.style={rounded corners=3pt, draw=#1, line width=1.1pt, fill=#1!8,
             text=#1!50!black, text width=2.45cm, align=flush left,
             inner sep=5pt, minimum height=2.5cm},
  hd/.style={rounded corners=2pt, fill=#1, text=white,
             font=\tiny\bfseries, inner sep=2.5pt, text width=2.3cm,
             align=center},
  a/.style={-{Stealth[length=2mm]}, neutral!70, line width=1pt}]

\node[st=primary]   (s1) at (0,0)    {\\[2pt]\texttt{"restart the}\\
  \texttt{server"}\\[4pt] split on whitespace and punctuation, fold case};
\node[st=secondary] (s2) at (3.1,0)  {\\[2pt]\texttt{restart}\\ \texttt{the}\\
  \texttt{server}\\[4pt] reduce inflected forms: \texttt{restarting}
  $\rightarrow$ \texttt{restart}};
\node[st=goldhl]    (s3) at (6.2,0)  {\\[2pt]\texttt{restart}/V\\
  \texttt{the}/Det\\ \texttt{server}/N\\[4pt] look up each word's possible
  categories};
\node[st=greenhl]   (s4) at (9.3,0)  {\\[2pt]a \emph{parse tree}\\[4pt]
  which words group with which, and how};
\node[st=accent]    (s5) at (12.4,0) {\\[2pt]\texttt{action=restart}\\
  \texttt{target=server}\\[4pt] something the router can act on};

\node[hd=primary]   at (s1.north) {1. TOKENISE};
\node[hd=secondary] at (s2.north) {2. MORPHOLOGY};
\node[hd=goldhl]    at (s3.north) {3. TAG};
\node[hd=greenhl]   at (s4.north) {4. PARSE};
\node[hd=accent]    at (s5.north) {5. INTERPRET};

\foreach \x/\y in {s1/s2,s2/s3,s3/s4,s4/s5}{\draw[a] (\x) -- (\y);}

\node[rounded corners=3pt, fill=lightbg, draw=primary!30, line width=0.8pt,
      text=primary!70!black, font=\scriptsize, text width=13.6cm,
      align=flush left, inner sep=6pt] at (6.2,-2.2)
  {\textbf{Each stage can fail, and a failure at stage 4 is the useful kind.} If the
   parser returns nothing, the subject line was not in the sub-language the grammar
   describes, and the system knows it does not know. If it returns \emph{two} trees,
   the request is genuinely ambiguous and a person should be asked. Both are
   better outcomes than a confident guess, and both come free with a grammar.};
\end{tikzpicture}}%
{The symbolic pipeline. Stages 1 to 3 are bookkeeping; stage 4 is where the
knowledge lives.}{nlp-pipeline}

\section{Context-Free Grammars}

\begin{definitionbox}[Context-free grammar]
A \term{context-free grammar} is a set of rules $A \to \beta$, where $A$ is a
\term{non-terminal} (a category such as \texttt{NP}) and $\beta$ a sequence of
non-terminals and \term{terminals} (words). A string is \term{in the language} if
the start symbol can be rewritten to it.

``Context-free'' means a rule may be applied wherever its left side appears,
regardless of surroundings --- which is what makes efficient parsing possible and
what makes the formalism too weak for some natural-language phenomena.
\end{definitionbox}

The grammar for ticket subjects is seven rules:

\smallskip
\noindent\begin{tabular}{@{}L{5.2cm}L{9.6cm}@{}}
\toprule
\textbf{Rule} & \textbf{Reading} \\
\midrule
$S \to \mathit{VP}$ & an imperative: ``restart the server'' \\
$S \to \mathit{NP}\ \mathit{VP}$ & a statement: ``the server crashed'' \\
$\mathit{VP} \to V\ \mathit{NP}$ & a verb and its object \\
$\mathit{VP} \to \mathit{VP}\ \mathit{PP}$ & a verb phrase modified by a
phrase \\
$\mathit{NP} \to \mathit{Det}\ N$ & a determiner and a noun \\
$\mathit{NP} \to \mathit{NP}\ \mathit{PP}$ & a noun phrase modified by a
phrase \\
$\mathit{PP} \to P\ \mathit{NP}$ & a preposition and its object \\
$N \to N\ N$ & a noun compound: ``payments gateway'' \\
\bottomrule
\end{tabular}

\smallskip
Every rule has either two non-terminals or one terminal on its right, which is
\term{Chomsky normal form} --- the shape the CYK algorithm needs.

\section{Parsing with CYK}

\term{CYK} is dynamic programming over spans. Fill a triangular chart: cell
$(i,j)$ holds every category that could cover words $i$ to $j$. Single words come
from the lexicon; longer spans are built by trying every split point $k$ and every
rule $A \to B\,C$, entering $A$ whenever $B$ covers $(i,k)$ and $C$ covers
$(k,j)$. The sentence parses if the start symbol reaches the top cell.

Its cost is $O(n^{3}|G|)$ --- cubic in the sentence length --- which for a subject
line of seven words means a chart of $28$ cells, of which $18$ are non-empty. The
important property is that it finds \emph{all} parses, and it finds them without
backtracking.

\section{Ambiguity Is the Point}

\begin{examplebox}[Five subject lines, parsed]
\begin{tabular}{@{}L{8.7cm}L{2.0cm}L{3.0cm}@{}}
\toprule
\textbf{Subject} & \textbf{Parses?} & \textbf{How many} \\
\midrule
\texttt{restart the server} & yes & 1 \\
\texttt{check the payments gateway} & yes & 1 \\
\texttt{restart the server after the timeout} & yes & \textbf{2} \\
\texttt{fix the login error on the gateway} & yes & \textbf{2} \\
\texttt{restart the gateway after the payments error} & yes & \textbf{2} \\
\midrule
\texttt{the the server restart} & no & 0 \\
\bottomrule
\end{tabular}
\end{examplebox}

Three of the five are ambiguous, and the ambiguity is not a quirk of the grammar.
Consider:

\begin{quote}\ttfamily\small
restart the server after the timeout
\end{quote}

Does \emph{after the timeout} attach to the verb phrase --- restart it, and do so
after the timeout occurs --- or to the noun phrase, meaning \emph{the server that
handles timeouts}? The two readings give different instructions, and the grammar
correctly reports that the sentence supports both.

\dsfig{%
\begin{tikzpicture}[font=\scriptsize,
  nt/.style={font=\scriptsize\bfseries, text=#1, inner sep=1.5pt},
  wd/.style={font=\ttfamily\tiny, text=neutral!60!black, inner sep=1.5pt},
  e/.style={neutral!60, line width=0.7pt}]

% ---- reading 1: PP attaches to VP
\node[nt=primary] (s1) at (0,3.0) {S};
\node[nt=primary] (vp1) at (0,2.3) {VP};
\node[nt=greenhl] (vpa) at (-1.5,1.6) {VP};
\node[nt=accent]  (pp1) at (1.8,1.6) {PP};
\node[nt=greenhl] (v1) at (-2.6,0.9) {V};
\node[nt=greenhl] (np1) at (-0.7,0.9) {NP};
\node[nt=accent]  (p1) at (0.9,0.9) {P};
\node[nt=accent]  (np2) at (2.7,0.9) {NP};
\node[wd] (w1) at (-2.6,0.25) {restart};
\node[wd] (w2) at (-1.3,0.25) {the};
\node[wd] (w3) at (-0.2,0.25) {server};
\node[wd] (w4) at (0.9,0.25) {after};
\node[wd] (w5) at (2.1,0.25) {the};
\node[wd] (w6) at (3.3,0.25) {timeout};
\foreach \a/\b in {s1/vp1,vp1/vpa,vp1/pp1,vpa/v1,vpa/np1,pp1/p1,pp1/np2,
                   v1/w1,np1/w2,np1/w3,p1/w4,np2/w5,np2/w6}{
  \draw[e] (\a) -- (\b);}
\node[font=\scriptsize\bfseries, text=greenhl!50!black] at (0,3.7)
  {Reading 1: attach to the \emph{verb}};
\node[font=\tiny\itshape, text=greenhl!50!black, text width=4.4cm,
      align=center] at (0,-0.35)
  {\emph{restart it, and do so after the timeout happens}};

% ---- reading 2: PP attaches to NP
\node[nt=primary] (s2) at (8.6,3.0) {S};
\node[nt=primary] (vp2) at (8.6,2.3) {VP};
\node[nt=greenhl] (v2) at (6.6,1.6) {V};
\node[nt=accent]  (npb) at (9.6,1.6) {NP};
\node[nt=accent]  (npc) at (8.3,0.9) {NP};
\node[nt=accent]  (pp2) at (10.9,0.9) {PP};
\node[wd] (x1) at (6.6,0.25) {restart};
\node[wd] (x2) at (7.7,0.25) {the};
\node[wd] (x3) at (8.8,0.25) {server};
\node[wd] (x4) at (9.9,0.25) {after};
\node[wd] (x5) at (11.0,0.25) {the};
\node[wd] (x6) at (12.1,0.25) {timeout};
\foreach \a/\b in {s2/vp2,vp2/v2,vp2/npb,npb/npc,npb/pp2,
                   v2/x1,npc/x2,npc/x3,pp2/x4,pp2/x5,pp2/x6}{
  \draw[e] (\a) -- (\b);}
\node[font=\scriptsize\bfseries, text=accent!55!black] at (8.6,3.7)
  {Reading 2: attach to the \emph{noun}};
\node[font=\tiny\itshape, text=accent!55!black, text width=4.4cm,
      align=center] at (9.3,-0.35)
  {\emph{restart the server that handles timeouts}};

\node[rounded corners=3pt, fill=lightbg, draw=primary!30, line width=0.9pt,
      text=primary!70!black, font=\scriptsize, text width=13.4cm,
      align=flush left, inner sep=6pt] at (4.8,-1.55)
  {\textbf{Two different instructions from one subject line.} A parser that
   returned only the first would be confidently wrong half the time and would give
   no indication of it. Returning both is the honest answer, and it tells the
   routing system exactly when to ask a person --- which is a capability the
   grammar provides for free.};
\end{tikzpicture}}%
{Prepositional-phrase attachment ambiguity. The words and their order are identical;
only the structure differs, and the structure is the meaning.}{parse-trees}

Ambiguity multiplies. Each additional prepositional phrase roughly doubles the
number of parses --- the counts follow the Catalan numbers --- so a sentence with
six of them has hundreds of readings. Resolving them requires knowledge the grammar
does not contain, which is where statistical and, later, neural methods enter.

\section{Pattern Substitution: Eliza}
\label{sec:eliza}

In 1966 Joseph Weizenbaum wrote a program of a few hundred lines that conducted
conversations in the manner of a Rogerian psychotherapist. It had no grammar, no
parser, no representation of the conversation and no model of anything. It matched
patterns and rearranged pronouns.

\begin{examplebox}[A service-desk Eliza]
\begin{tabular}{@{}L{6.6cm}L{8.2cm}@{}}
\toprule
\textbf{User says} & \textbf{Rule that matched, and the reply} \\
\midrule
\texttt{I am worried about the release} &
\texttt{I am *} $\to$ \texttt{How long have you been *?}\newline
``How long have you been worried about the release?'' \\
\texttt{my deploy keeps failing} &
\texttt{my * keeps *} $\to$ \texttt{Tell me more about your *.}\newline
``Tell me more about your deploy.'' \\
\texttt{because the config is wrong} &
\texttt{because *} $\to$ \texttt{Is that the real reason?}\newline
``Is that the real reason?'' \\
\texttt{no idea what is happening} &
(no pattern) $\to$ fallback\newline
``Can you tell me more about that?'' \\
\bottomrule
\end{tabular}
\end{examplebox}

The mechanism is exactly \chref{fol}'s unification, restricted to strings: a pattern
containing a variable is matched against an input, the variable binds to whatever
it covered, and the binding is substituted into a template. Weizenbaum added one
more trick --- swap \emph{my} for \emph{your} and \emph{I} for \emph{you} in the
captured text --- and that is the whole program.

\begin{alertbox}[The Eliza effect]
Weizenbaum's own account is that people who knew exactly how the program worked ---
his secretary among them --- nonetheless asked to be left alone with it, and
disclosed things they would not have told a colleague. He found this so disturbing
that he spent much of the rest of his career arguing against the uses being proposed
for such systems.

The lesson named after the episode --- the \term{Eliza effect} --- is that
\emph{people attribute understanding to any system that produces fluent,
contextually apt responses}, and that the attribution is nearly impossible to
suppress even with full knowledge of the mechanism.

That is worth holding onto when you meet \chref{trends}. The question is not whether
a system's output feels like understanding; Eliza's did, in two hundred lines, in
1966. The question is what the system \emph{represents} and what it can be held to
--- which is what \chref{introduction}'s five-point checklist asks, and why the
checklist is the right instrument rather than your impression of the transcript.
\end{alertbox}

\section{N-Grams and Their Limit}

The alternative to grammar is counting. An \term{n-gram} model estimates the
probability of a word from the $n-1$ words before it:
\[
  P(w_{i} \mid w_{1} \ldots w_{i-1}) \;\approx\;
  P(w_{i} \mid w_{i-n+1} \ldots w_{i-1})
\]
This is the \term{Markov assumption}: everything before the window is irrelevant.
Estimating the right-hand side is just counting, which is why n-gram models were
the backbone of speech recognition and machine translation for decades and why they
still make excellent autocompletions for a constrained domain such as ticket
subjects.

Two limitations, and only one of them is fixable.

\textbf{Sparsity} is fixable. Most possible trigrams never appear, even in a large
corpus, so unseen ones get probability zero and the model assigns zero to any
sentence containing one. Smoothing --- reserving probability mass for the unseen ---
handles it, and \emph{Machine Learning} Chapter~10 covers the standard methods.

\textbf{The fixed window is not.} In
\emph{the server that the engineer who wrote the runbook maintained was restarted},
the subject of \emph{was restarted} is \emph{the server}, nine words back. No
$n$-gram with a practical $n$ can connect them, and increasing $n$ makes sparsity
exponentially worse. The dependency is a structural fact about the sentence, and
capturing it is exactly what a grammar does and what counting cannot.

That tension --- grammars capture structure but are brittle, counting is robust but
structurally blind --- defined the field for thirty years. \chref{trends} describes
what resolved it.

\begin{worked}[Filling the CYK chart, by hand]
Parse \texttt{restart the server} with the grammar above. Three words, so the chart
is a triangle of six cells.

\smallskip
\textbf{Row 1 --- single words, from the lexicon.}
\begin{tight}
  \item $(0,1)$ \texttt{restart}: $\{V\}$.
  \item $(1,2)$ \texttt{the}: $\{\mathit{Det}\}$.
  \item $(2,3)$ \texttt{server}: $\{N\}$.
\end{tight}

\smallskip
\textbf{Row 2 --- spans of two.}
\begin{tight}
  \item $(0,2)$ covering \texttt{restart the}: the only split is $k=1$, giving
        $V$ then $\mathit{Det}$. No rule has right side $V\ \mathit{Det}$, so the
        cell is \textbf{empty}.
  \item $(1,3)$ covering \texttt{the server}: split at $k=2$ gives
        $\mathit{Det}$ then $N$, and $\mathit{NP} \to \mathit{Det}\ N$ applies.
        Cell: $\{\mathit{NP}\}$.
\end{tight}

\smallskip
\textbf{Row 3 --- the whole sentence, $(0,3)$.} Two split points:
\begin{tight}
  \item $k=1$: cell $(0,1) = \{V\}$ and cell $(1,3) = \{\mathit{NP}\}$. The rule
        $\mathit{VP} \to V\ \mathit{NP}$ applies, so $\mathit{VP}$ is entered.
  \item $k=2$: cell $(0,2)$ is empty, so nothing.
\end{tight}
Then the unit rule $S \to \mathit{VP}$ adds $S$. Cell $(0,3) = \{\mathit{VP}, S\}$,
the start symbol is present, and the sentence parses --- by exactly one derivation.

\smallskip
\textbf{Now the ambiguous case.} Add \texttt{after the timeout}. At the top cell
there are now \emph{two} ways to reach $S$: one via
$\mathit{VP} \to \mathit{VP}\ \mathit{PP}$ (split after \texttt{server}), and one
via $\mathit{VP} \to V\ \mathit{NP}$ where the $\mathit{NP}$ was itself built by
$\mathit{NP} \to \mathit{NP}\ \mathit{PP}$ (split after \texttt{restart}). Both
entries record $S$ in the same cell, which is why counting parses means counting
\emph{derivations} rather than checking membership --- and why the lab keeps
back-pointers.
\end{worked}

\begin{pitfall}
\begin{tight}
  \item \textbf{Treating a single parse as the answer.} Three of the five subject
        lines here have two readings. A parser that returns the first is
        confidently wrong and silent about it.
  \item \textbf{Forgetting Chomsky normal form.} CYK needs rules of the form
        $A \to B\,C$ or $A \to w$. A grammar written naturally must be converted
        first.
  \item \textbf{Assuming a bigger grammar is a better one.} Every added rule
        multiplies the ambiguity. Real broad-coverage grammars return thousands of
        parses per sentence, which is why ranking them became the whole problem.
  \item \textbf{Mistaking Eliza's fluency for understanding.} It has no
        representation of anything. If the impression of understanding is your
        evidence, you have no evidence.
  \item \textbf{Believing more data fixes the n-gram window.} Sparsity is fixable
        by smoothing; a fixed window cannot represent a nine-word dependency at any
        corpus size.
  \item \textbf{Parsing when you could match.} If the domain is genuinely a few
        dozen command shapes, patterns are shorter, faster and easier to maintain.
        Reach for a grammar when structure matters.
\end{tight}
\end{pitfall}

%---------------------------------------------------------------------
\begin{lab}[Parsing Ticket Subjects]
A CYK parser with back-pointers, the ambiguity count, an Eliza responder, and a
bigram model with its limitation demonstrated.
\begin{lstlisting}[language=Python]
"""Chapter 18 lab -- symbolic language processing on ticket subjects.

A CYK chart parser, a pattern-substitution responder in the style of
Eliza, and a bigram model. Nothing here learns; the knowledge is in
the grammar and the patterns.
"""

import re
from collections import defaultdict

# --- the grammar, in Chomsky normal form ---------------------------
RULES = [
    ("S", ("VP",)),                 # imperative: "restart the server"
    ("S", ("NP", "VP")),            # statement:  "the server crashed"
    ("VP", ("V", "NP")),
    ("VP", ("VP", "PP")),           # PP attaches to the verb phrase
    ("NP", ("Det", "N")),
    ("NP", ("NP", "PP")),           # ... or to the noun phrase
    ("PP", ("P", "NP")),
    ("N", ("N", "N")),              # noun compound: "payments gateway"
]
LEXICON = {
    "restart": ["V"], "check": ["V"], "fix": ["V"],
    "server": ["N"], "payments": ["N"], "login": ["N"],
    "error": ["N"], "timeout": ["N"], "gateway": ["N"],
    "the": ["Det"], "a": ["Det"],
    "on": ["P"], "after": ["P"], "in": ["P"],
}


def cyk(words):
    """Fill the triangular chart. Returns (chart, back-pointers)."""
    n = len(words)
    chart = defaultdict(set)
    back = defaultdict(list)

    def close(cell):
        """Apply unit rules such as S -> VP until nothing changes."""
        changed = True
        while changed:
            changed = False
            for (lhs, rhs) in RULES:
                if len(rhs) == 1 and rhs[0] in chart[cell] \
                        and lhs not in chart[cell]:
                    chart[cell].add(lhs)
                    changed = True

    for i, w in enumerate(words):
        for tag in LEXICON.get(w, []):
            chart[(i, i + 1)].add(tag)
        close((i, i + 1))

    for span in range(2, n + 1):
        for i in range(0, n - span + 1):
            j = i + span
            for k in range(i + 1, j):
                for (lhs, rhs) in RULES:
                    if len(rhs) != 2:
                        continue
                    b, c = rhs
                    if b in chart[(i, k)] and c in chart[(k, j)]:
                        chart[(i, j)].add(lhs)
                        back[(i, j, lhs)].append((k, b, c))
            close((i, j))
    return chart, back


def count_parses(chart, back, i, j, sym, memo=None):
    """Count DERIVATIONS, not membership: an ambiguous sentence
    reaches the same cell by two different routes."""
    if memo is None:
        memo = {}
    key = (i, j, sym)
    if key in memo:
        return memo[key]
    total = 0
    for (k, b, c) in back[key]:
        total += (count_parses(chart, back, i, k, b, memo)
                  * count_parses(chart, back, k, j, c, memo))
    for (lhs, rhs) in RULES:                     # unit rules
        if lhs == sym and len(rhs) == 1 and rhs[0] != sym \
                and rhs[0] in chart[(i, j)]:
            total += count_parses(chart, back, i, j, rhs[0], memo)
    if total == 0 and j == i + 1 and sym in chart[(i, j)]:
        total = 1
    memo[key] = total
    return total


def parse(subject):
    words = subject.split()
    chart, back = cyk(words)
    ok = "S" in chart[(0, len(words))]
    n = count_parses(chart, back, 0, len(words), "S") if ok else 0
    return ok, n, chart


# --- Eliza: pattern substitution, and nothing else ------------------
REFLECT = {"my": "your", "i": "you", "me": "you", "am": "are",
           "mine": "yours", "myself": "yourself"}

PATTERNS = [
    (r"^i am (.*)$", "How long have you been {0}?"),
    (r"^my (\w+) keeps (.*)$", "Tell me more about your {0}."),
    (r"^because (.*)$", "Is that the real reason?"),
    (r"^i (?:can't|cannot) (.*)$", "What stops you from {0}?"),
    (r"^(.*) is broken$", "What makes you say {0} is broken?"),
]
FALLBACK = "Can you tell me more about that?"


def reflect(text):
    return " ".join(REFLECT.get(w, w) for w in text.split())


def eliza(line):
    low = line.lower().strip().rstrip(".!?")
    for (pat, template) in PATTERNS:
        m = re.match(pat, low)
        if m:
            return template.format(*[reflect(g) for g in m.groups()])
    return FALLBACK


# --- a bigram model -------------------------------------------------
CORPUS = [
    "restart the server", "restart the gateway", "restart the service",
    "check the payments gateway", "check the login error",
    "check the server", "fix the login error", "fix the timeout error",
    "restart the payments gateway", "check the timeout",
]


def bigrams(corpus):
    counts = defaultdict(lambda: defaultdict(int))
    for line in corpus:
        toks = ["<s>"] + line.split()
        for a, b in zip(toks, toks[1:]):
            counts[a][b] += 1
    return counts


def suggest(counts, word, k=3):
    nxt = counts.get(word, {})
    total = sum(nxt.values())
    ranked = sorted(nxt.items(), key=lambda kv: -kv[1])[:k]
    return [(w, c / total) for w, c in ranked]


SUBJECTS = [
    "restart the server",
    "check the payments gateway",
    "restart the server after the timeout",
    "fix the login error on the gateway",
    "restart the gateway after the payments error",
]

if __name__ == "__main__":
    print("=== parsing ticket subjects ===")
    print("subject                                         parses  count")
    print("-" * 62)
    counts = {}
    for s in SUBJECTS:
        ok, n, _ = parse(s)
        counts[s] = n
        print("%-46s %7s %6d" % (s, ok, n))
    # three of the five are structurally ambiguous
    ambiguous = [s for s in SUBJECTS if counts[s] > 1]
    assert len(ambiguous) == 3 and all(counts[s] == 2 for s in ambiguous)
    print()
    print("%d of %d subjects have more than one parse."
          % (len(ambiguous), len(SUBJECTS)))
    print("Each extra prepositional phrase roughly doubles the count.")

    print()
    print("=== and one that is not in the language ===")
    for bad in ("the the server restart", "server the",
                "restart restart the"):
        ok, n, _ = parse(bad)
        print("   %-32s parses: %s" % (bad, ok))
        assert not ok
    print("   Returning nothing is a useful answer: the system knows")
    print("   that it does not know.")

    print()
    print("=== the chart ===")
    words = SUBJECTS[4].split()
    chart, _ = cyk(words)
    cells = sum(1 for k in chart if chart[k])
    total = len(words) * (len(words) + 1) // 2
    print("   %d words -> %d cells, %d of them non-empty"
          % (len(words), total, cells))

    print()
    print("=== Eliza ===")
    for line in ("I am worried about the release",
                 "my deploy keeps failing",
                 "because the config is wrong",
                 "I can't reproduce it",
                 "the gateway is broken",
                 "no idea what is happening"):
        print("   > %-34s %s" % (line, eliza(line)))
    print()
    print("   Pattern matching and pronoun swapping. No grammar, no")
    print("   representation of the conversation, no model of")
    print("   anything -- and it reads as though it is listening.")
    assert eliza("something unmatched entirely") == FALLBACK

    print()
    print("=== a bigram model ===")
    bg = bigrams(CORPUS)
    for w in ("<s>", "restart", "the", "login"):
        print("   after %-10s -> %s"
              % (w, ", ".join("%s %.2f" % t for t in suggest(bg, w))))

    print()
    print("=== the limitation no corpus size removes ===")
    hard = ("the server that the engineer who wrote the runbook "
            "maintained was restarted")
    toks = hard.split()
    print("   '%s'" % hard)
    print("   'was restarted' is at position %d; its subject 'server'"
          % toks.index("was"))
    print("   is at position %d -- %d words back."
          % (toks.index("server"), toks.index("was") - toks.index("server")))
    print("   No n-gram with a practical n connects them, and raising")
    print("   n makes sparsity exponentially worse. A grammar has no")
    print("   difficulty with it, because the dependency is structural.")
    assert toks.index("was") - toks.index("server") >= 8
\end{lstlisting}
Then add one more prepositional phrase --- \texttt{restart the server after the
timeout on the gateway} --- and count the parses. The answer is five, not four: the
growth follows the Catalan numbers, and this is the moment broad-coverage parsing
stops being a matter of writing more rules.
\end{lab}

%---------------------------------------------------------------------
\begin{checkpoint}
\begin{tightnum}
  \item \texttt{restart the server after the timeout} has exactly two parses. Give
        both readings in English, and say what a routing system should do with a
        subject line that has more than one.
  \item Eliza has no grammar, no parser and no representation of the conversation,
        and people disclosed things to it they would not tell a colleague. State
        what this shows about fluency as evidence of understanding.
  \item Smoothing fixes n-gram sparsity but nothing fixes the fixed window. Explain
        the difference, with the nine-word dependency from this chapter.
\end{tightnum}
\end{checkpoint}

%---------------------------------------------------------------------
\begin{chaptersummary}
\begin{tight}
  \item The symbolic pipeline runs tokenise, morphology, tag, parse, interpret.
        Stages 1--3 are bookkeeping; the \emph{parse} is where the knowledge lives.
  \item A \term{context-free grammar} is rules $A \to \beta$. In \term{Chomsky
        normal form} every rule is $A \to B\,C$ or $A \to w$, which is what
        \term{CYK} needs.
  \item CYK is dynamic programming over spans, costs $O(n^{3}|G|)$, and finds
        \emph{every} parse without backtracking.
  \item \term{Ambiguity} is the normal case, not a defect: three of five subject
        lines here have two readings, and each extra prepositional phrase roughly
        doubles the count. Returning two trees is the honest answer and tells the
        system when to ask a person.
  \item \term{Eliza} matched patterns and swapped pronouns in a few hundred lines,
        and people attributed understanding to it. The \term{Eliza effect} is that
        fluency compels that attribution regardless of what is underneath --- which
        is why \chref{introduction}'s checklist, and not your impression of a
        transcript, is the instrument to use.
  \item Pattern substitution is \chref{fol}'s unification restricted to strings.
  \item An \term{n-gram} model applies the \term{Markov assumption} and is pure
        counting. \term{Sparsity} is fixable by smoothing; the \emph{fixed window}
        is not, and a nine-word dependency defeats it at any corpus size. Grammars
        capture structure and are brittle; counting is robust and structurally
        blind.
\end{tight}
\end{chaptersummary}

\begin{reviewq}
  \item \textbf{(MCQ)} CYK requires the grammar to be in: (a)~Backus--Naur form
        (b)~Chomsky normal form (c)~Greibach normal form (d)~any context-free
        form.
  \item \textbf{(MCQ)} The time complexity of CYK in the sentence length is:
        (a)~$O(n)$ (b)~$O(n^{2})$ (c)~$O(n^{3})$ (d)~$O(2^{n})$.
  \item \textbf{(MCQ)} Eliza's mechanism is closest to: (a)~chart parsing
        (b)~unification restricted to strings (c)~an n-gram model (d)~forward
        chaining.
  \item \textbf{(MCQ)} Smoothing an n-gram model addresses: (a)~the fixed window
        (b)~sparsity (c)~structural ambiguity (d)~the Markov assumption.
  \item \textbf{(Short)} Give the two readings of \texttt{fix the login error on the
        gateway} and say which one a routing system would need to disambiguate
        before acting.
  \item \textbf{(Short)} Explain why returning no parse is a more useful failure
        than returning a guess, with a service-desk example.
  \item \textbf{(Short)} State the Eliza effect and say why it is relevant to
        evaluating a modern language system.
  \item \textbf{(Applied)} Add the rule $\mathit{VP} \to V$ (an intransitive verb,
        as in \texttt{restart}) to the grammar, and re-count the parses of
        \texttt{restart the server after the timeout}. Explain the change, and say
        what it illustrates about adding rules to a grammar.
\end{reviewq}
